% GENERATED FILE. Edit canonical lessons, then run npm run generate:books.
% canonical-source-hash: fnv1a64:6d65753c630b49d1
% canonical-lessons: JA-C392-yorimichi, JA-C392-mawarimichi, JA-C392-ikidomari, JA-C392-hikikaesu, JA-C392-tachiyoru

\chapter{Stopping on the Way}
\label{ch:ja-stopping-on-the-way}
\hlchaptermodality{392}

\begin{chapteropening}
\textbf{By the end of this chapter:} \emph{I can talk about detours, dead ends and stops along the way.}

Complete the last lesson of chapter 392: I can talk about detours, dead ends and stops along the way.
\end{chapteropening}

\section[\texorpdfstring{\ja{よりみち} (yorimichi)}{よりみち (yorimichi)}]{\ja{よりみち} (yorimichi) — a stop on the way; a side trip}

\label{lesson:JA-C392-yorimichi}



\begin{quote}
Before the new one: say the Japanese for to stop walking, then the Japanese for to block.
\end{quote}

\subsection*{You'll want to know: \ja{よりみち}}
\textbf{\ja{よりみち}} — \emph{yorimichi} — \textquotedblleft{}a stop on the way; a side trip\textquotedblright{}.

\begin{quote}
\textbf{\ja{しごとの} \ja{あとで、ほんやに} \ja{よりみちしました。}} — \emph{shigoto no ato de, hon'ya ni yorimichi shimashita} — \textquotedblleft{}After work I stopped off at a bookshop.\textquotedblright{}
\end{quote}

\begin{cousinweb}
\textbf{\ja{よりみち}} (\emph{yorimichi}) is a native Japanese word (wago), inherited rather than borrowed from Chinese, built from \emph{yori}, \textquotedblleft{}dropping in\textquotedblright{} (from \emph{yoru}) and \emph{michi}, \textquotedblleft{}road\textquotedblright{}; its kanji were attached later as kun'yomi (JMdict).
\end{cousinweb}

\begin{grammarlens}[title={What you've built}]
One more word for getting around.
\end{grammarlens}

\subsection*{Your turn}
\begin{itemize}
\raggedright
  \item \emph{Say it:} \emph{yorimichi}
  \item \emph{Say it:} \emph{yorimichi}, once more
  \item \emph{Say it:} \emph{shigoto no ato de, hon'ya ni yorimichi shimashita} — After work I stopped off at a bookshop.
  \item \emph{Recall:} say the Japanese for to stop walking, then the Japanese for to block, then say \emph{yorimichi} again
\end{itemize}

\begin{tcolorbox}[breakable,colback=teal!4,colframe=teal!35!black,title={Before you move on}]
What does \ja{よりみち} mean? (\textquotedblleft{}A stop on the way; a side trip\textquotedblright{}.) Say it once more.
\end{tcolorbox}
\section[\texorpdfstring{\ja{まわりみち} (mawarimichi)}{まわりみち (mawarimichi)}]{\ja{まわりみち} (mawarimichi) — a detour; a roundabout way}

\label{lesson:JA-C392-mawarimichi}



\begin{quote}
Before the new one: say the Japanese for to block, then the Japanese for a stop on the way.
\end{quote}

\subsection*{You'll want to know: \ja{まわりみち}}
\textbf{\ja{まわりみち}} — \emph{mawarimichi} — \textquotedblleft{}a detour; a roundabout way\textquotedblright{}.

\begin{quote}
\textbf{\ja{じゅうたいで} \ja{まわりみちを} \ja{しました。}} — \emph{jūtai de mawarimichi o shimashita} — \textquotedblleft{}Because of the traffic jam we took a detour.\textquotedblright{}
\end{quote}

\begin{cousinweb}
\textbf{\ja{まわりみち}} (\emph{mawarimichi}) is a native Japanese word (wago), inherited rather than borrowed from Chinese, built from \emph{mawari}, \textquotedblleft{}going round\textquotedblright{} (from \emph{mawaru}) and \emph{michi}, \textquotedblleft{}road\textquotedblright{}; its kanji were attached later as kun'yomi (JMdict).
\end{cousinweb}

\begin{grammarlens}[title={What you've built}]
One more word for getting around.
\end{grammarlens}

\subsection*{Your turn}
\begin{itemize}
\raggedright
  \item \emph{Say it:} \emph{mawarimichi}
  \item \emph{Say it:} \emph{mawarimichi}, once more
  \item \emph{Say it:} \emph{jūtai de mawarimichi o shimashita} — Because of the traffic jam we took a detour.
  \item \emph{Recall:} say the Japanese for to block, then the Japanese for a stop on the way, then say \emph{mawarimichi} again
\end{itemize}

\begin{tcolorbox}[breakable,colback=teal!4,colframe=teal!35!black,title={Before you move on}]
What does \ja{まわりみち} mean? (\textquotedblleft{}A detour; a roundabout way\textquotedblright{}.) Say it once more.
\end{tcolorbox}
\section[\texorpdfstring{\ja{いきどまり} (ikidomari)}{いきどまり (ikidomari)}]{\ja{いきどまり} (ikidomari) — a dead end}

\label{lesson:JA-C392-ikidomari}



\begin{quote}
Before the new one: say the Japanese for a stop on the way, then the Japanese for a detour.
\end{quote}

\subsection*{You'll want to know: \ja{いきどまり}}
\textbf{\ja{いきどまり}} — \emph{ikidomari} — \textquotedblleft{}a dead end\textquotedblright{}.

\begin{quote}
\textbf{\ja{この} \ja{みちは} \ja{いきどまりです。}} — \emph{kono michi wa ikidomari desu} — \textquotedblleft{}This road is a dead end.\textquotedblright{}
\end{quote}

\begin{cousinweb}
\textbf{\ja{いきどまり}} (\emph{ikidomari}) is a native Japanese word (wago), inherited rather than borrowed from Chinese, built from \emph{iki}, \textquotedblleft{}going\textquotedblright{} (from \emph{iku}) and \emph{domari}, \textquotedblleft{}stopping\textquotedblright{} (from \emph{tomaru}); its kanji were attached later as kun'yomi (JMdict).
\end{cousinweb}

\begin{grammarlens}[title={What you've built}]
One more word for getting around.
\end{grammarlens}

\subsection*{Your turn}
\begin{itemize}
\raggedright
  \item \emph{Say it:} \emph{ikidomari}
  \item \emph{Say it:} \emph{ikidomari}, once more
  \item \emph{Say it:} \emph{kono michi wa ikidomari desu} — This road is a dead end.
  \item \emph{Recall:} say the Japanese for a stop on the way, then the Japanese for a detour, then say \emph{ikidomari} again
\end{itemize}

\begin{tcolorbox}[breakable,colback=teal!4,colframe=teal!35!black,title={Before you move on}]
What does \ja{いきどまり} mean? (\textquotedblleft{}A dead end\textquotedblright{}.) Say it once more.
\end{tcolorbox}
\section[\texorpdfstring{\ja{ひきかえす} (hikikaesu)}{ひきかえす (hikikaesu)}]{\ja{ひきかえす} (hikikaesu) — to turn back}

\label{lesson:JA-C392-hikikaesu}



\begin{quote}
Before the new one: say the Japanese for a detour, then the Japanese for a dead end.
\end{quote}

\subsection*{You'll want to know: \ja{ひきかえす}}
\textbf{\ja{ひきかえす}} — \emph{hikikaesu} — \textquotedblleft{}to turn back\textquotedblright{}.

\begin{quote}
\textbf{\ja{みちを} \ja{まちがえたので、ひきかえしました。}} — \emph{michi o machigaeta node, hikikaeshimashita} — \textquotedblleft{}I took the wrong road, so I turned back.\textquotedblright{}
\end{quote}

Polite \textbf{\ja{ひきかえします}}, \emph{hikikaeshimasu}; a godan verb. Past (polite): \textbf{\ja{ちちが} \ja{えきまで} \ja{ひきかえしました。}}, \emph{chichi ga eki made hikikaeshimashita}.

\begin{cousinweb}
\textbf{\ja{ひきかえす}} (\emph{hikikaesu}) is a native Japanese word (wago), inherited rather than borrowed from Chinese, built from \emph{hiki}, \textquotedblleft{}pulling\textquotedblright{} (from \emph{hiku}) and \emph{kaesu}, \textquotedblleft{}to send back\textquotedblright{}; its kanji were attached later as kun'yomi (JMdict).
\end{cousinweb}

\begin{grammarlens}[title={What you've built}]
One more word for getting around.
\end{grammarlens}

\subsection*{Your turn}
\begin{itemize}
\raggedright
  \item \emph{Say it:} \emph{hikikaesu}
  \item \emph{Say it:} \emph{hikikaesu}, once more
  \item \emph{Say it:} \emph{michi o machigaeta node, hikikaeshimashita} — I took the wrong road, so I turned back.
  \item \emph{Recall:} say the Japanese for a detour, then the Japanese for a dead end, then say \emph{hikikaesu} again
\end{itemize}

\begin{tcolorbox}[breakable,colback=teal!4,colframe=teal!35!black,title={Before you move on}]
What does \ja{ひきかえす} mean? (\textquotedblleft{}To turn back\textquotedblright{}.) Say it once more.
\end{tcolorbox}
\section[\texorpdfstring{\ja{たちよる} (tachiyoru)}{たちよる (tachiyoru)}]{\ja{たちよる} (tachiyoru) — to stop by; to drop in on the way}

\label{lesson:JA-C392-tachiyoru}



\begin{quote}
Before the new one: say the Japanese for a dead end, then the Japanese for to turn back.
\end{quote}

\subsection*{You'll want to know: \ja{たちよる}}
\textbf{\ja{たちよる}} — \emph{tachiyoru} — \textquotedblleft{}to stop by; to drop in on the way\textquotedblright{}.

\begin{quote}
\textbf{\ja{えきの} \ja{まえの} \ja{きっさてんに} \ja{たちよりました。}} — \emph{eki no mae no kissaten ni tachiyorimashita} — \textquotedblleft{}I stopped by a cafe in front of the station.\textquotedblright{}
\end{quote}

Polite \textbf{\ja{たちよります}}, \emph{tachiyorimasu}; a godan verb. Past (polite): \textbf{\ja{あねが} \ja{はなやに} \ja{たちよりました。}}, \emph{ane ga hanaya ni tachiyorimashita}.

\begin{cousinweb}
\textbf{\ja{たちよる}} (\emph{tachiyoru}) is a native Japanese word (wago), inherited rather than borrowed from Chinese, built from \emph{tachi}, \textquotedblleft{}standing\textquotedblright{} (from \emph{tatsu}) and \emph{yoru}, \textquotedblleft{}to drop in\textquotedblright{}; its kanji were attached later as kun'yomi (JMdict).
\end{cousinweb}

\begin{grammarlens}[title={What you've built}]
One more word for getting around.
\end{grammarlens}

\subsection*{Your turn}
\begin{itemize}
\raggedright
  \item \emph{Say it:} \emph{tachiyoru}
  \item \emph{Say it:} \emph{tachiyoru}, once more
  \item \emph{Say it:} \emph{eki no mae no kissaten ni tachiyorimashita} — I stopped by a cafe in front of the station.
  \item \emph{Recall:} say the Japanese for a dead end, then the Japanese for to turn back, then say \emph{tachiyoru} again
\end{itemize}

\begin{tcolorbox}[breakable,colback=teal!4,colframe=teal!35!black,title={Before you move on}]
What does \ja{たちよる} mean? (\textquotedblleft{}To stop by; to drop in on the way\textquotedblright{}.) Say it once more.
\end{tcolorbox}
